%% =============================================================================
%% BLOCO DE SUBSTITUIÇÃO — Cap. 4, Metodologia.
%% Substitui TODA a antiga "\section{Segunda Campanha: Verificação sobre o
%% Código Original}" (sec:segunda-campanha) e suas subseções, até imediatamente
%% ANTES de "\section{Resumo}".
%%
%% MOTIVO (12/11/2026): a versão anterior descrevia uma "campanha de saída" de
%% dez módulos (pca9685, dshot, voxl_esc, tap_esc, uavcan, px4io, mixer_module,
%% control_allocation, control_allocator) com artefatos em "saida/". Esse
%% diretório NÃO existe no repositório e metade dos recortes estava "a definir":
%% a campanha não foi executada. A segunda campanha REAL, com log, é a destas
%% três frentes (candidatos_extras_v5 + RoboClaw + batmon), coerente com a
%% introdução e a capa. Todo número abaixo vem de log do repositório.
%% =============================================================================

% =============================================================================
\section{\corr{Segunda Campanha: Código de Produção, \textit{Replay} e Sanitizadores}}
\label{sec:segunda-campanha}
% =============================================================================

\corr{A primeira campanha mostrou a limitação de verificar trechos transcritos:
o veredito vale para a cópia, e a conclusão sobre o driver depende de uma
conferência que, nos casos do FIFO e do GPS, inverteu o resultado. A segunda
campanha foi desenhada para reduzir essa distância. Ela mantém as oito etapas,
mas endurece a conferência contra o código de produção (etapa~v): cada candidato
alcançável é submetido a \textit{replay} do \emph{código real} do PX4 sob
sanitizadores dinâmicos, e a validação cruzada com o CBMC passa a cobrir todos os
\textit{harnesses}, não apenas os candidatos. Nenhum achado desta campanha foi
reportado à comunidade: qualquer \textit{report} depende da revisão do Prof.\
Lucas.}

\corr{A campanha reúne dez módulos complementares aos doze da primeira, em três
frentes: (A)~oito candidatos adicionais em módulos de entrada e de sistema;
(B)~o driver de saída \textbf{RoboClaw}; e (C)~o driver de bateria
\textbf{batmon}. As frentes compartilham o protocolo de evidência --- verificação
com ESBMC e CBMC, \textit{replay} do código real e sanitizadores ---, mas
diferem no commit e no desfecho, como se detalha a seguir.}

\subsection{\corr{Frente A --- Oito candidatos adicionais}}
\label{subsec:2c-candidatos}

\corr{Oito módulos foram verificados no mesmo commit da primeira campanha
(\texttt{8be2590c15}), com ESBMC~8.1.0 e~8.5.0, CBMC~5.95.1 e os sanitizadores de
endereço (ASan) e de comportamento indefinido (UBSan). Um único script
(\texttt{rodar\_tudo.sh}) regenera todos os registros em qualquer máquina com as
três ferramentas, cada log com cabeçalho (data, máquina, pasta, comando e
versões) e código de saída diferente de zero se algum resultado divergir do
esperado. A Tabela~\ref{tab:candidatos-extras} resume os desfechos.}

\begin{table}[!htb]
\centering
\caption{\corr{Candidatos adicionais da segunda campanha (commit \texttt{8be2590c15})}}
\label{tab:candidatos-extras}
\scriptsize
\renewcommand{\arraystretch}{0.95}
\setlength{\tabcolsep}{3pt}
\begin{tabular}{|p{2.4cm}|p{5.2cm}|p{5.0cm}|}
\hline
\textbf{Módulo} & \textbf{Resultado no commit testado} & \textbf{Situação e desfecho} \\
\hline
\corr{\texttt{sumd.cpp}} & \corr{com quadro legítimo de 32 canais, \texttt{sumd\_decode()} escreve o 64.º byte em \texttt{sumd\_data[64]}, sobre \texttt{crc16\_high}} & \corr{\textbf{continua no \texttt{main}}; defeito reproduzido por \textit{replay} da função real; \textit{issue} redigida, pendente de revisão} \\
\hline
\corr{\texttt{MicroStrain.cpp}} & \corr{construtor: \texttt{max\_len} pode ser 128 e \texttt{\_port[128]='\textbackslash0'} escreve fora do vetor de 128} & \corr{\textbf{continua no \texttt{main}}; defeito reproduzido (ASan); \textit{issue} redigida, pendente de revisão} \\
\hline
\corr{\texttt{externalChecks.cpp}} & \corr{\texttt{num\_events} (0--255) indexa \texttt{events[5]}} & \corr{era real; \textbf{corrigido \textit{upstream}} em \texttt{ca1d36476a} (PR~\#28598)} \\
\hline
\corr{\texttt{ft7\_technologies.cpp}} & \corr{contadores de byte indexam \texttt{\_readbuf[30]} sem limite} & \corr{era real; \textbf{corrigido \textit{upstream}} em \texttt{3acf7935a0} (PR~\#28600)} \\
\hline
\corr{\texttt{IridiumSBD.cpp}} & \corr{a guarda de espaço compara com \texttt{\_packet\_length}, mas o \texttt{memcpy} copia \texttt{buflen}} & \corr{era real; \textbf{corrigido \textit{upstream}} em \texttt{0c903d47d6} (PR~\#28599)} \\
\hline
\corr{\texttt{bst.cpp}} & \corr{violação apenas no \textit{harness} que omitia duas guardas já presentes no commit (de \texttt{e8e86a2e0f})} & \corr{\textbf{falso positivo de modelagem}; o \textit{harness} fiel verifica \textsc{successful}} \\
\hline
\corr{\texttt{logger.cpp}} & \corr{nenhuma violação de memória em \texttt{write\_info\_multiple()} para qualquer \texttt{key\_len}} & \corr{sem defeito de memória; o truncamento é alcançável, mas gera registro com \texttt{key\_len} errado, não escrita fora dos limites} \\
\hline
\corr{\texttt{Common.cpp}} & \corr{violação só com índice que a guarda de produção nunca permite} & \corr{\textbf{retirado}} \\
\hline
\end{tabular}
\end{table}

\corr{Dois pontos merecem registro. Primeiro, três candidatos
(\texttt{externalChecks}, \texttt{ft7} e \texttt{IridiumSBD}) eram violações reais
no commit testado e foram \emph{corrigidos pelo próprio PX4}, por outras pessoas,
em setembro de 2026; as correções \textit{upstream} corroboram, de forma
independente, que os achados eram defeitos reais, e por isso nenhuma \textit{issue}
foi aberta para eles. Segundo, a verificação com o CBMC usou as mesmas
pré-condições e limites de desdobramento do ESBMC, com \textit{unwinding
assertions} e \textit{bounds} que cobrem os laços (por exemplo, 260 no
\texttt{externalChecks} e 345 no \texttt{IridiumSBD}); as duas ferramentas
acusam as mesmas violações. O caso do \texttt{sumd} precisou de
\texttt{--unwind 70}; como o modo \texttt{--multi-property} do ESBMC ultrapassa
quinze minutos nesse limite, a contagem de propriedades foi tomada do CBMC.}

\subsection{\corr{Frente B --- Driver de saída RoboClaw}}
\label{subsec:2c-roboclaw}

\corr{O driver \texttt{roboclaw} converte os comandos de controle em pacotes do
protocolo \textit{Packet Serial} enviados por UART aos controladores de motor.
Dois defeitos foram verificados no commit \texttt{8be2590c15}, com ESBMC~8.2.0
(solver Z3) e CBMC~5.95.1 (Tabela~\ref{tab:roboclaw-resultados}).}

\begin{table}[!htb]
\centering
\caption{\corr{RoboClaw: verificação formal dos dois defeitos}}
\label{tab:roboclaw-resultados}
\footnotesize
\begin{tabular}{|p{7.6cm}|p{2.6cm}|p{2.6cm}|}
\hline
\textbf{Defeito} & \textbf{Sem correção} & \textbf{Com \texttt{-DCORRECAO}} \\
\hline
\corr{Bug 1 --- byte de velocidade $>127$ em \texttt{sendUnsigned7Bit()}} & \corr{\textsc{failed}} & \corr{\textsc{successful}} \\
\hline
\corr{Bug 2 --- erro de \texttt{write()} não detectado em \texttt{writeCommandWithPayload()}} & \corr{\textsc{failed}} & \corr{\textsc{successful}} \\
\hline
\end{tabular}

\scriptsize
\begin{flushleft}
\corr{Mesmo padrão nas duas ferramentas (ESBMC e CBMC), na máquina da autora, em
28/10/2026. Registros e \textit{hashes} em
\texttt{testes\_hardware/roboclaw/evidencia\_maquina/}.}
\end{flushleft}
\end{table}

\corr{O \textbf{Bug~1} decorre do caminho de terminação de voo: o
\texttt{module.yaml} admite valores de \textit{failsafe} até 257 e o
\textit{mixer} os repassa ao driver sem limitar, enquanto o manual da RoboClaw
(Rev.~6.0, pp.~48--49) restringe o byte de dados dos comandos 0, 1, 4 e 5 à faixa
$0$--$127$. Um valor de \textit{failsafe} igual a 257 produz o byte 130. O
\textit{replay} com UART emulada, sobre a lógica real do driver, envia de fato o
byte 130 de forma determinística. O \textbf{Bug~2} vem de armazenar o retorno de
\texttt{write()} num \texttt{size\_t}: quando \texttt{write()} devolve $-1$, o valor
vira um inteiro sem sinal grande, a checagem \texttt{bytes\_written < 2} não
detecta o erro e a função retorna \texttt{OK}; o driver então espera uma resposta
que não virá, e a falha só aparece como \textit{timeout} em \texttt{select()}.
A correção dos dois está em \texttt{correcao\_roboclaw.patch}.}

\corr{\paragraph{Validação em hardware do Bug~1.} O Bug~1 foi observado também
em bancada, numa Pixhawk~6X~RT com firmware compilado do mesmo commit
(\texttt{px4\_fmu-v6xrt\_default}). A porta TELEM2 foi ligada a um ESP32 que
faz o papel da RoboClaw: responde às consultas do driver e registra cada byte
recebido. Com \texttt{RBCLW\_FAIL1} igual a 257 e a terminação de voo acionada
pelo console (\texttt{commander termination}), a placa passou a transmitir o
pacote \texttt{80 00 82 8A 90}: endereço \texttt{0x80}, comando~0, byte de
dados 130 e CRC válido. O pacote se repetiu 24\,136 vezes nas duas capturas
guardadas, sem erro de CRC. Dois ensaios complementares, na mesma montagem,
isolam a causa (Tabela~\ref{tab:roboclaw-hardware}): com o parâmetro em 254, o
maior valor que ainda resulta em byte válido, o mesmo firmware envia 127; e,
com a correção aplicada e o parâmetro de volta em 257, o byte também fica em
127, em 5\,267 pacotes e nenhum acima da faixa.}

\begin{table}[htb]
\centering
\caption{\corr{RoboClaw: byte de dados observado na serial da Pixhawk~6X~RT durante a terminação}}
\label{tab:roboclaw-hardware}
\begin{tabular}{llll}
\hline
\corr{\textbf{Firmware}} & \corr{\textbf{\texttt{RBCLW\_FAIL1}}} & \corr{\textbf{Pacote}} & \corr{\textbf{Byte de dados}} \\
\hline
\corr{Original} & \corr{257} & \corr{\texttt{80 00 82 8A 90}} & \corr{130 (fora de 0--127)} \\
\corr{Original} & \corr{254} & \corr{\texttt{80 00 7F B4 22}} & \corr{127} \\
\corr{Corrigido} & \corr{257} & \corr{\texttt{80 00 7F B4 22}} & \corr{127} \\
\hline
\end{tabular}
\end{table}

\corr{Esse resultado tem quatro limites, que delimitam o que ele prova.
Primeiro, a RoboClaw é emulada: o ensaio mostra que a placa \emph{transmite} o
byte inválido, mas não o que o controlador real faz ao recebê-lo, o que o
manual não especifica. Segundo, o firmware não é o de fábrica: o driver não
integra a configuração padrão dessa placa e foi incluído com
\texttt{CONFIG\_DRIVERS\_ROBOCLAW=y}; além disso, foi preciso inicializar dois
membros da classe (\texttt{\_uart\_fd\_set} e \texttt{\_uart\_fd\_timeout}),
sem o que o driver bloqueava a fila de trabalho na partida. Essa alteração não
toca o caminho de saída. Terceiro, a placa registrou \textit{timeout} de ACK
durante os ensaios, de modo que o driver foi observado no caminho de erro de
resposta; a medida, que é o byte lido diretamente do fio, não depende disso.
Quarto, só o Bug~1 foi exercitado em hardware; a evidência do Bug~2 continua
sendo a formal e o \textit{replay} com UART emulada. As capturas, os
\textit{diffs} dos firmwares e o código do emulador estão em
\texttt{testes\_hardware/roboclaw/hardware\_pixhawk6xrt/}.}

\subsection{\corr{Frente C --- Driver de bateria batmon}}
\label{subsec:2c-batmon}

\corr{No driver \texttt{smart\_battery/batmon}, o número de células lido do
registrador SMBus da bateria define \texttt{\_cell\_count} até
\texttt{MAX\_CELL\_COUNT}${}=16$, e o laço de publicação escreve
\texttt{new\_report.voltage\_cell\_v[i]} para \texttt{i < \_cell\_count}. Mas
\texttt{voltage\_cell\_v} é um \texttt{float32[14]}
(\texttt{msg/versioned/BatteryStatus.msg}). Uma bateria legítima de 15S ou 16S,
dentro dos registradores que o próprio driver define, já basta para escrever fora
do vetor. Com 16 células, o valor extra cai em \texttt{full\_charge\_capacity\_wh}
(o \textit{replay} publica 3,715, a tensão da 16.ª célula), e \texttt{cell\_count}${}=16$
é publicado num vetor de 14 posições. O driver está nas configurações padrão de
várias placas (fmu-v3 a v6u, cubeorange, entre outras).}

\corr{A Tabela~\ref{tab:batmon-resultados} reúne os resultados. Diferentemente das
outras frentes, a conferência foi feita contra o \texttt{main} de 13/10/2026
(commit \texttt{c9b9e4e25e}), e não \texttt{8be2590c15} --- ressalva que se
mantém explícita.}

\begin{table}[!htb]
\centering
\caption{\corr{batmon: resultados (commit \texttt{c9b9e4e25e})}}
\label{tab:batmon-resultados}
\footnotesize
\begin{tabular}{|p{6.8cm}|p{3.2cm}|p{2.6cm}|}
\hline
\textbf{Ferramenta / verificação} & \textbf{Código atual} & \textbf{\texttt{-DCORRECAO}} \\
\hline
\corr{ESBMC~8.1.0 (\texttt{--unwind 20 --overflow-check})} & \corr{\textsc{failed}} & \corr{\textsc{successful}} \\
\hline
\corr{ESBMC~8.5.0 (\texttt{--multi-property})} & \corr{\textsc{failed} (2 de 13: \texttt{i < 14} e \textit{array bounds})} & \corr{\textsc{successful}} \\
\hline
\corr{CBMC~5.95.1 (\texttt{--unwind 20}, \textit{bounds} e \textit{pointer check})} & \corr{1 de 15 falha} & \corr{0 de 15} \\
\hline
\corr{UBSan (\texttt{./replay 16})} & \corr{\texttt{index 14 out of bounds for type 'float [14]'}} & \corr{---} \\
\hline
\corr{\textit{replay} sem sanitizador} & \corr{\texttt{full\_charge\_capacity\_wh=3.7150}} & \corr{---} \\
\hline
\end{tabular}
\end{table}

\corr{O mesmo defeito foi corrigido no driver UAVCAN/DroneCAN
(\texttt{cbat\_sub\_cb}) pelo PR~\#28606, em 14/09/2026, relatado por um terceiro;
o \texttt{batmon} ficou de fora dessa correção, e as buscas por ``batmon'',
``cell\_count'' e ``voltage\_cell\_v'' nas \textit{issues} do PX4 não encontraram
registro. O achado está, portanto, em aberto, e seu \textit{report} depende da
revisão do Prof.\ Lucas.}

\subsection{\corr{Protocolo de conferência contra o código de produção}}
\label{subsec:2c-protocolo}

\corr{As três frentes seguem o mesmo protocolo de evidência, que operacionaliza a
etapa~(v) (Seção~\ref{subsec:etapa-replay}). As entradas do código verificado são
declaradas não determinísticas e restritas por pré-condições
(\texttt{\_\_VERIFIER\_assume}, aceita pelo ESBMC e pelo CBMC), com a origem de
cada pré-condição anotada no cabeçalho do \textit{harness}: um limite de
parâmetro declarado no \texttt{module.yaml}, o modo como o próprio PX4 chama a
função, ou uma constante do código. No RoboClaw, por exemplo, a saída recebida do
\textit{mixer} foi restrita à faixa $[0;\,257]$ porque o \texttt{module.yaml}
declara \texttt{failsafe: \{min: 0, max: 257\}}; uma pré-condição mais estreita
teria ocultado o Bug~1.}

\corr{Cada candidato alcançável passa por \textit{replay}: a função \emph{real}
do PX4 é compilada com um compilador comum (g++), executada com os valores do
contraexemplo e observada sob ASan e UBSan. O critério é binário --- se o código
real reproduz o efeito previsto (o byte 130 no RoboClaw, o acesso a
\texttt{voltage\_cell\_v[14]} no batmon, a escrita em \texttt{\_port[128]} no
MicroStrain), o candidato é defeito reproduzido; se não reproduz, investiga-se a
guarda de produção que o \textit{harness} omitiu, como no \texttt{bst}. Quando há
correção, o mesmo \textit{harness} é reverificado com o \textit{patch}
(\texttt{-DCORRECAO}) e deve passar; para os três candidatos corrigidos
\textit{upstream}, a reverificação usa a correção do próprio PX4
(\texttt{-DCORRECAO\_UPSTREAM}).}

\corr{A validação cruzada com o CBMC foi estendida a todos os \textit{harnesses},
com as opções mapeadas para os equivalentes mais próximos
(Tabela~\ref{tab:equivalencia-opcoes}). Divergências entre as ferramentas são
registradas como resultado da comparação: elas vêm de modelagem de biblioteca
(por exemplo, \texttt{isfinite} tratado pelo CBMC como retorno arbitrário, que
exigiu um prelúdio de modelagem) ou de diferença de versão na contagem de
propriedades, e não de desacordo sobre a existência das violações.}

\begin{table}[!htb]
\centering
\caption{\corr{Equivalência de opções entre ESBMC e CBMC~5.95}}
\label{tab:equivalencia-opcoes}
\footnotesize
\begin{tabular}{|p{6.5cm}|p{7.0cm}|}
\hline
\textbf{ESBMC} & \textbf{CBMC} \\
\hline
\corr{checagens padrão (limites, ponteiros, divisão por zero)} & \corr{\texttt{--bounds-check --pointer-check --div-by-zero-check}} \\
\hline
\corr{\texttt{--overflow-check}} & \corr{\texttt{--signed-overflow-check --float-overflow-check}} \\
\hline
\corr{com \textit{unwinding assertions} (padrão) / \texttt{--no-unwinding-assertions}} & \corr{\texttt{--unwinding-assertions} / padrão} \\
\hline
\corr{\textit{solver} Boolector ou Z3} & \corr{MiniSat (padrão)} \\
\hline
\end{tabular}
\end{table}

\corr{\paragraph{Classificação dos desfechos.} A segunda campanha usa a mesma
escala do Capítulo~\ref{cap:introducao}, com uma distinção que a primeira já
exigia: um \textbf{falso positivo de modelagem} decorre de uma guarda de produção
omitida no \textit{harness} (o \texttt{bst} aqui, como o contador de FIFO na
primeira campanha), ao passo que um \textbf{falso positivo da ferramenta} decorre
do verificador, não do código (a modelagem de \texttt{isfinite} pelo CBMC). Os
desfechos da campanha distribuem-se assim: \textbf{defeito reproduzido no código
de produção} (\texttt{sumd}, \texttt{MicroStrain}, RoboClaw, batmon);
\textbf{defeito real corroborado por correção \textit{upstream}}
(\texttt{externalChecks}, \texttt{ft7}, \texttt{IridiumSBD}); \textbf{falso
positivo de modelagem} (\texttt{bst}); e \textbf{sem violação} (\texttt{logger},
além do \texttt{Common} retirado).}

\pendente{Inserir os números das \textit{issues} do \texttt{sumd} e do
\texttt{MicroStrain}, e das eventuais \textit{issues} do RoboClaw e do batmon,
após a revisão do Prof.\ Lucas e o envio.}
